\section*{Apêndice}
\addcontentsline{toc}{section}{Apêndice}

\subsection*{Levantamento Preliminar de Literatura}
Todo o Levantamento Preliminar de Literatura foi realizado no Banco de Teses e Dissertações da CAPES no período em 18/12/2025, na área do conhecimento de Ciências da Computação, com o recorte temporal de 2013 a 2025. Foram utilizadas palavras-chave relacionadas aos temas centrais do projeto combinadas com operadores booleanos.
\subsubsection{Visão Computacional}

\begin{table}[H]
\centering
\caption{Levantamento preliminar de trabalhos relacionados à visão computacional}
\renewcommand{\arraystretch}{1.4}
\label{tab:lev_prel_vc}

% 1. Diminui o espaço entre as colunas para ganhar área de texto
\setlength{\tabcolsep}{3pt} 
\setstretch{1.1}
\scriptsize

\begin{tabularx}{\textwidth}{|
>{\RaggedRight\arraybackslash}p{1.4cm}| % Autor
>{\centering\arraybackslash}p{0.6cm}|    % T/D
>{\RaggedRight\arraybackslash}X|        % Objetivo
>{\RaggedRight\arraybackslash}X|        % Contexto
>{\RaggedRight\arraybackslash}X|        % Metodologia
>{\RaggedRight\arraybackslash}p{2.3cm}|        % Resultados
>{\RaggedRight\arraybackslash}X|        % Limitações
>{\RaggedRight\arraybackslash}p{2.1cm}| % Contribuições (Tamanho fixo evita estouro)
}
\hline
\textbf{Autor} & \textbf{T/D} & \textbf{Objetivo} & \textbf{Contexto} & \textbf{Metodologia} & \textbf{Resultados} & \textbf{Limitações} & \textbf{Contribuições} \\ 

\hline
\citeonline{Cardoso2020} & 
D & 
Automatizar a contagem de moscas em bovinos. & 
Tecnologias computacionais para agropecuária. & 
Experimental. & 
Desenvolvimento de um método computacional eficaz. & 
Base de imagens pequena. & 
Reforça as aplicações da VC. \\ 

\hline
\citeonline{Cunha2022} & 
D & 
Implementar rastreamento ocular de baixo custo. & 
Ambiente educacional. & 
Experimental. & 
Desenvolvimento de sistema embarcado eficiente. & 
Precisão inferior a soluções comerciais. & 
Mostra a viabilidade de soluções de VC. \\ 

\hline
\citeonline{Menezes2013}  & 
D & 
Analisar abordagens que melhorem o desempenho da servovisão. & 
Robótica aérea. & 
Experimental. & 
Aponta vantagens de controladores alternativos. & 
Avaliação restrita a cenários específicos. & 
Demonstra alternativas para o controle de robôs utilizando VC \\ 

\hline
\end{tabularx}

\vspace{2mm}
\end{table}

\subsubsection{Internet das Coisas (IOT) e Aprendizado Profundo (DL)}

\begin{table}[H]
\centering
\caption{Levantamento preliminar de trabalhos relacionados à Internet das Coisas (IOT) e Aprendizado Profundo (DL)}
\renewcommand{\arraystretch}{1.4}
\label{tab:lev_prel_iot_dl}

% 1. Diminui o espaço entre as colunas para ganhar área de texto
\setlength{\tabcolsep}{3pt} 
\setstretch{1.1}
\scriptsize

\begin{tabularx}{\textwidth}{|
>{\RaggedRight\arraybackslash}p{1.4cm}| % Autor
>{\centering\arraybackslash}p{0.6cm}|    % T/D
>{\RaggedRight\arraybackslash}X|        % Objetivo
>{\RaggedRight\arraybackslash}X|        % Contexto
>{\RaggedRight\arraybackslash}X|        % Metodologia
>{\RaggedRight\arraybackslash}p{2.3cm}|        % Resultados
>{\RaggedRight\arraybackslash}X|        % Limitações
>{\RaggedRight\arraybackslash}p{2.1cm}| % Contribuições (Tamanho fixo evita estouro)
}
\hline
\textbf{Autor} & \textbf{T/D} & \textbf{Objetivo} & \textbf{Contexto} & \textbf{Metodologia} & \textbf{Resultados} & \textbf{Limitações} & \textbf{Contribuições} \\ 

\hline
\citeonline{Pereira2022} & 
D & 
Propor um método de classificação de tráfego de rede com DL. & 
Redes definidas por software integradas a IoT. & 
Experimental. & 
Alta eficiência na identificação de aplicações em ambiente SDN & 
Dependência de dados estatísticos específicos do OpenFlow. & 
Mostra a aplicabilidade de Deep Learning em SDN para IoT e Indústria 4.0. \\ 

\hline
\citeonline{GuimaraesDaSilva2023} & 
T & 
Consumo de Energia em sistemas IoT & 
Avaliar modelos de DL. & 
Experimental. & 
Apresentou uma comparação entre os principais modelos.  & 
Avaliação restrita a séries temporais univariadas. & 
Fornece insights sobre a eficácia de diferentes modelos de DL em IoT. \\ 

\hline
\citeonline{BritoGuimaraes2023}  & 
D & 
Avaliar modelos para detecção de botnets dispositivos IoT. & 
Segurança cibernética em Internet das Coisas. & 
Experimental. & 
Aponta modelos mais eficazes para detecção de botnets em ambientes IoT reais. & 
Dependência de hardware de borda relativamente potente (Jetson Nano). & 
Introduz a dimensão de segurança inteligente em sistemas robóticos. \\ 

\hline
\end{tabularx}

\vspace{2mm}
\end{table}

\subsubsection{}

\begin{table}[H]
\centering
\caption{Levantamento preliminar de trabalhos relacionados à Internet das Coisas (IOT) e Aprendizado Profundo (DL)}
\renewcommand{\arraystretch}{1.4}
\label{tab:lev_prel_iot_dl}

% 1. Diminui o espaço entre as colunas para ganhar área de texto
\setlength{\tabcolsep}{3pt} 
\setstretch{1.1}
\scriptsize

\begin{tabularx}{\textwidth}{|
>{\RaggedRight\arraybackslash}p{1.4cm}| % Autor
>{\centering\arraybackslash}p{0.6cm}|    % T/D
>{\RaggedRight\arraybackslash}X|        % Objetivo
>{\RaggedRight\arraybackslash}X|        % Contexto
>{\RaggedRight\arraybackslash}X|        % Metodologia
>{\RaggedRight\arraybackslash}p{2.3cm}|        % Resultados
>{\RaggedRight\arraybackslash}X|        % Limitações
>{\RaggedRight\arraybackslash}p{2.1cm}| % Contribuições (Tamanho fixo evita estouro)
}
\hline
\textbf{Autor} & \textbf{T/D} & \textbf{Objetivo} & \textbf{Contexto} & \textbf{Metodologia} & \textbf{Resultados} & \textbf{Limitações} & \textbf{Contribuições} \\ 

\hline
\citeonline{Rezeck2019} & 
D & 
Projetar e validar uma plataforma de Robótica educacional & 
Robótica educacional. & 
Experimental. & 
Plataforma funcional de baixo custo. & 
Capacidade computacional e sensorial restrita devido ao foco em baixo custo. & 
Demonstra viabilidade de robôs educacionais abertos e acessíveis. \\ 

\hline
\citeonline{Bertoncini2022} & 
T & 
Desenvolver uma central de controle para navegação. & 
Robótica móvel autônoma aplicada. & 
Experimental. & 
Sistema capaz de detectar obstáculos e comandar desvios com comportamento satisfatório.  & 
Sensível à iluminação e reflexos, exigindo calibração prévia do ambiente. & 
Fornece uma base conceitual para navegação visual centralizada. \\ 

\hline
\citeonline{Santos2022}  & 
D & 
Construir um manipulador robótico baseado em VC. & 
Automação industrial. & 
Experimental. & 
Precisão nas tarefas de pick-and-place.  & 
Dependência de melhorias mecânicas e eletromecânicas. & 
Fornece uma base para integração de visão computacional, IA e controle de manipuladores. \\ 

\hline
\end{tabularx}

\vspace{2mm}
\end{table}